% Section II — Objectives
% Standalone prose form for department template (v2.2 STRUCT-02)

The scientific objective of this thesis was to build and evaluate a
local-first Vietnamese financial phishing detection system that keeps a
user's messages private while still explaining its answers well enough to
act on. Underneath, the task is supervised four-class text classification:
every message has one correct label, and the Qwen path also produces a risk
tier, the suspicious phrases quoted from the message, and a short written
explanation. The finished system adds safe next steps on top of those, and
checks each one against a fixed safety rule (Section~\ref{sec:problem-statement}).

The objective has seven parts, as follows:

\begin{enumerate}
  \item Build a Vietnamese phishing dataset from real scam-warning
  articles. This project checks every message automatically for three
  things: that all the required fields are present, that the label is one
  of the four allowed classes, and that the highlighted suspicious phrases
  can be found in the message text. This project writes several messages
  from each source article, and two messages from the same article can
  differ by little more than a bank name or an amount, so this project puts
  all of them in the same set. No article ever has some of its messages
  used for training and others kept for testing.

  \item Review part of the dataset by hand instead of trusting the judge
  model. This project draws the sample evenly across all four classes and
  across both judge verdicts, so it contains the messages the judge passed
  as well as the ones it failed, not only the easy ones. Where the reviewer
  disagreed with a label, this project applied the correction only if those
  messages did not all trace back to one source article. A large group of
  near-identical messages is not a large group of examples, and accepting
  them together would make one class look far more varied than it really
  is.

  \item Fix the messages that came out badly. A review found that one whole
  class had been written as third-person narration instead of as real chat
  messages, and that other classes carried leftover template text. This
  project rewrote those messages offline from their original scenarios.

  \item Choose the base model by measuring it rather than by reputation.
  This project ran three Qwen models over the same 33 example messages and
  compared them on how many risky messages each one caught, how much video
  memory it needed, and how long it took per message.

  \item Check whether fine-tuning was needed at all. This project scored
  the chosen model on the same messages twice, once with prompting alone
  and once after fine-tuning, so that training the model rested on a
  measurement rather than on the assumption that training must help.

  \item Measure ordinary LoRA and QLoRA on the actual RTX~5050 laptop, and
  choose between them from the recorded time and memory numbers rather than
  from an assumption that one of them would run out of memory.

  \item Train Qwen QLoRA and PhoBERT locally, export the chosen Qwen model
  as a Q8\_0 GGUF file and confirm it loads, then evaluate both finished
  models exactly once on the same 220 held-out messages. This project
  separated those messages from the rest at the start, and neither model,
  nor any decision taken during the project, was allowed to see them until
  the training was over.
\end{enumerate}

This project met most of these objectives. It built and split the dataset,
chose the base model and the training method from measurements, trained
both models on the laptop, exported the Qwen model, and finished the final
evaluation on both without producing a single unreadable answer. The
numbers are in Chapter~V.

This project left two further pieces of work unfinished, on purpose in one
case. It did not fit a deployment version of the model after the
evaluation, because the point of a single final evaluation is lost if
training continues on the same data afterwards. It also reorganised the
code to be easier to read and navigate, but a later review found security
problems that are still open, so that part is not complete either.
